\section{Model and the readiness objective}
\label{sec:model}

Let \(N=\{1,\ldots,n\}\) be a finite collection of independent modules.
Module \(i\) has preparation capacity \(M_i>0\) and releases source capacity
\(a_i\geq0\) when it transfers. The initial spare source capacity is \(s\geq0\).
Preparation \(p_i\) belongs to \([0,M_i]\). Before a removal request, no
module transfers, and nonnegative preparation allocations \(v_i\) and
optional work \(u\) satisfy
\begin{equation}
 u+\sum_{i\in N}v_i\leq s.
 \label{eq:normal-budget}
\end{equation}
Following a request, optional work stops. If \(D\) is the transferred set,
the available preparation capacity is
\begin{equation}
 b(D)=s+\sum_{j\in D}a_j.
 \label{eq:released-capacity}
\end{equation}
A module can transfer when its preparation equals \(M_i\). Transfer is
instantaneous, removes that module's preparation obligation, and releases
\(a_i\) immediately. Releases and transfers have no cross-module constraints.

The main structural results concern proportional losses
\(g_i(p)=\gamma_i p\), with \(\gamma_i>0\). The scheduling reduction will be
proved for the larger class
\begin{equation}
 g_i:[0,M_i]\longrightarrow[0,\infty)
 \quad\text{nondecreasing and Lipschitz},\qquad g_i(0)=0.
 \label{eq:smooth-class}
\end{equation}
Allowed actual losses are measurable and obey
\(0\leq\rho_i(t)\leq g_i(p_i(t))\).
The uncertainty set includes every simultaneous maximum feedback
\(\rho_i=g_i(p_i)\); it is a product constraint on the actual states.
Under maximum feedback, before transfer the dynamics are
\begin{equation}
 \dot p_i=
 \begin{cases}
 v_i-g_i(p_i),&0\leq p_i<M_i,\\
 \min\{v_i-g_i(M_i),0\},&p_i=M_i.
 \end{cases}
 \label{eq:reflected-dynamics}
\end{equation}
There is no lower reflection term because \(g_i(0)=0\) and \(v_i\geq0\).
The corresponding upper-reflected dynamics apply under smaller losses.

Policies are deterministic and causal, with measurable allocations. They may
prepare modules in parallel, interrupt preparation, and adapt to observations.
Every guarantee below quantifies over all allowed loss histories. In
particular, maximum feedback is an admissible adversary for every policy,
rather than a fixed numerical loss trace shared by different trajectories.
Scalar comparison implies that larger initial preparation and smaller losses
cannot worsen a feasible exit strategy. Transfers can likewise be performed
earlier without harming another module: capacity is nondecreasing and the
modules are independent.

Write \(F(p)\) for the infimum, over such policies starting at \(p\), of their
worst-case time to transfer all modules. Infinite values are allowed. For a
finite deadline \(H\geq0\), define the ready set and its minimum upkeep by
\begin{equation}
 \mathcal R_H=\{p\in\textstyle\prod_i[0,M_i]:F(p)\leq H\},
 \qquad
 C(H)=\min_{p\in\mathcal R_H}\sum_i g_i(p_i).
 \label{eq:readiness-cost}
\end{equation}
Attainment of this minimum is established below. A uniformly ready normal
policy must belong to \(\mathcal R_H\) at every counterfactual request time,
after every admissible normal loss history.

The model presupposes a sufficient transfer state and separately paid
receiver resources. It does not derive a proportional loss envelope from a
network or replication protocol. Initial preparation is paid before the
uniform-readiness guarantee begins unless a separate warmup construction is
given. Positive handoff latency, delayed capacity release, and interdependent
transfers are outside these assumptions.

All normal preparation input is charged to \eqref{eq:normal-budget}.
Accordingly, $C(H)$ measures maintenance through this specified source
bottleneck, excluding receiver provisioning and independent continuing
operation. The prohibition on normal-mode transfer is an additional service
requirement, not a consequence of preserving essential service alone.
Setting every $a_i=0$ gives an exact alternative interpretation: independent
perishable jobs must remain uncompleted until a command permits absorbing
completion. States, controls, losses and deadlines are unchanged. Thus
positive resource release and controller removal are not necessary for the
concentration property; its minimal interpretation is preparation for a
future completion command.

\section{Serial exit from arbitrary preparation}
\label{sec:seriality}

A full-capacity serial schedule transfers initially full modules at time zero
and then assigns the entire available capacity to one remaining module until
it transfers. Its order is chosen as part of the optimization.

\begin{lemma}[Postponing scalar input]
\label{lemma:postponement}
Extend a function \(g\) satisfying \eqref{eq:smooth-class} constantly beyond
\(M\). Consider the uncapped equation \(y'=v-g(y)\) with the same initial
condition under nonnegative bounded measurable inputs \(v_{\rm old}\) and
\(v_{\rm new}\). If
\[
 D(t)=\int_0^t(v_{\rm new}-v_{\rm old})(r)\,dr\leq0
 \quad(0\leq t\leq T),\qquad D(T)=0,
\]
then \(y_{\rm new}(T)\geq y_{\rm old}(T)\).
\end{lemma}
\begin{proof}
Both solutions remain nonnegative. Set \(z=y_{\rm new}-y_{\rm old}\) and
\[
 \alpha(t)=
 \begin{cases}
 \dfrac{g(y_{\rm new}(t))-g(y_{\rm old}(t))}{z(t)},&z(t)\ne0,\\
 0,&z(t)=0.
 \end{cases}
\]
Monotonicity and Lipschitz continuity give
\(0\leq\alpha(t)\leq\operatorname{Lip}(g)\) almost everywhere. Thus
\(z'=v_{\rm new}-v_{\rm old}-\alpha z\), and variation of constants gives
\[
 z(T)=\int_0^T k(t)D'(t)\,dt,\qquad
 k(t)=\exp\left(-\int_t^T\alpha(r)\,dr\right).
\]
Since \(k'=\alpha k\geq0\), integration by parts yields
\(z(T)=-\int_0^T k'(t)D(t)\,dt\geq0\).
\end{proof}

The uncapped equation is only a comparison device. In a constructed physical
schedule, transfer a module at its first hit of \(M\) and omit its subsequent
planned allocations. This increases available capacity. A module that has
not hit \(M\) agrees with its uncapped comparison trajectory.

\begin{theorem}[Arbitrary-state seriality]
\label{thm:seriality}
Under \eqref{eq:smooth-class}, every finite feasible maximum-loss exit schedule
can be replaced by a full-capacity serial schedule whose final completion is
no later. Consequently, if \(F(p)<\infty\), an optimizing serial order attains
the robust value \(F(p)\). The optimum is attained among finitely many orders.
No preservation of an arbitrary schedule's completion order is asserted.
\end{theorem}
\begin{proof}
Initially full modules can transfer immediately. Let \(b\) be the resulting
capacity, and consider the remaining subfull modules. We may likewise transfer
every module at its first hit of its target: delaying a transfer cannot help,
because its release is nonnegative and its remaining loss is then removed.
If \(b=0\), none can
increase, so no finite exit is possible. Otherwise take a finite schedule
under maximum loss, and let \(i\) be a first module to complete, at time \(T\).
Capacity is \(b\) before that time.

A subfull module can complete at constant capacity \(b\) only if
\(b>g_i(M_i)\). Indeed, if \(b\leq g_i(M_i)\), a root of \(b-g_i\) is an
equilibrium barrier for the full-rate scalar equation. Lipschitz uniqueness
prevents finite-time arrival from below or crossing of that root. Starting
above a lower root gives nonpositive full-rate drift and cannot reach the
larger target. Allocating less than \(b\) cannot improve this comparison.
Conversely, \(b>g_i(M_i)\) gives a strictly positive drift throughout the
preparation interval.

The time to prepare \(i\) immediately at rate \(b\) is therefore
\[
 \tau=\int_{p_i(0)}^{M_i}\frac{dq}{b-g_i(q)}.
\]
Let \(V_i=\int_0^T v_i(t)\,dt\) in the original schedule. For
\(W_i(p)=\int_0^p dq/(b-g_i(q))\), its trajectory satisfies
\[
 \frac{d}{dt}W_i(p_i(t))
 =\frac{v_i-g_i(p_i)}{b-g_i(p_i)}
 \leq\frac{v_i}{b},
\]
because \(0\leq v_i\leq b\) and \(g_i\geq0\). Integration gives
\(V_i\geq b\tau\). If \(W\) denotes all work allocated to the other modules
before \(T\), the shared budget gives
\begin{equation}
 W\leq b(T-\tau).
 \label{eq:first-stage-work}
\end{equation}

We now postpone all those other modules' work to \([\tau,T]\), retaining
each module's total allocation and never increasing its cumulative allocation
before \(T\). For \(j\ne i\), let
\(E_j=\int_0^\tau v_j(t)\,dt\). On \([\tau,T]\), retain its old allocation
and define spare rate
\[
 h(t)=b-\sum_{j\ne i}v_j(t)\geq0.
\]
Equation \eqref{eq:first-stage-work} implies
\[
 \int_\tau^T h(t)\,dt
 =b(T-\tau)-W+\sum_{j\ne i}E_j
 \geq\sum_{j\ne i}E_j.
\]
Order the other modules arbitrarily and write
\(B(t)=\int_\tau^t h(r)\,dr\). Cumulative repayment to module \(j\) is
\[
 A_j(t)=\min\left\{E_j,
       \max\left\{0,B(t)-\sum_{\ell\ {\rm before}\ j}E_\ell\right\}\right\}.
\]
These absolutely continuous functions satisfy \(A_j(T)=E_j\) and
\(\sum_j A_j'\leq h\) almost everywhere. Give other modules zero allocation
before \(\tau\), and allocation \(v_j+A_j'\) afterward. Their total rate is at
most \(b\). For each module the new cumulative input is no larger than the
old one at every intermediate time and is equal at \(T\).

Prepare \(i\) at rate \(b\) until \(\tau\) and transfer it. Its release makes
the postponed allocations feasible. Lemma~\ref{lemma:postponement} shows that
at \(T\) every other module has at least its original preparation, unless it
has already completed, which is better. Stop such newly completed modules
and omit their later planned allocations. Simultaneous original completions
cause no difficulty: all their targets have been reached by \(T\).

At \(T\) the constructed state and completed set dominate the originals.
Replay the original suffix while omitting allocations to already completed
modules. Scalar comparison and nondecreasing capacity preserve domination.
The resulting schedule finishes no later than the original and has a
full-capacity first stage. Apply the same argument to its residual schedule
at \(\tau\). Induction on the number of remaining modules gives a serial
schedule no slower than the original.

For a fixed initial vector there are finitely many serial orders. Every
finite stage has a unique full-rate completion time, and inaccessible orders
are assigned infinity. The least finite order value is therefore attained.
Any robust policy must face simultaneous maximum feedback, so its
worst-case time is at least this value. Conversely, execute an optimizing
order against a virtual maximum-loss trajectory. Under smaller losses,
transfer each module no later than its virtual completion; greater
preparation and earlier capacity release preserve the comparison for the
remaining stages. This policy attains the same bound for every allowed loss
history and proves the robust assertion.
\end{proof}

\begin{lemma}[Compact readiness]
\label{lemma:compactness}
For every finite \(H\geq0\), \(\mathcal R_H\) is nonempty and compact.
Consequently the minimum in \eqref{eq:readiness-cost} is attained.
\end{lemma}
\begin{proof}
The all-full state exits at time zero. To prove closedness without assuming
continuity of hitting times at rate barriers, fix a permutation and
nonnegative stage durations \(\ell_1,\ldots,\ell_n\) with
\(\sum_r\ell_r\leq H\). During stage \(r\), allocate the full capacity of its
completed prefix to the selected module; unserved modules decay passively.
Require every selected module to be full at its stage endpoint. A full
module may voluntarily wait for its specified stage, and a zero-duration
stage is allowed when it is already full.

Passive flows and fixed-capacity upper-reflected scalar flows are continuous
in initial state and elapsed time. This follows from the Lipschitz drift
estimate and scalar comparison; upper reflection does not increase the
distance between ordered scalar states. Each endpoint constraint is
therefore closed in the compact product of the preparation box and the
duration simplex. Its projection onto the initial states is compact.
Taking the finite union over permutations remains compact.

Every projected witness is a feasible schedule, including deliberate
delays of already full modules. Conversely, Theorem~\ref{thm:seriality}
supplies such a witness for every state in \(\mathcal R_H\). Thus this
projection equals the ready set, including initially full and inaccessible-rate
boundaries. The upkeep objective is continuous, so it attains its minimum.
\end{proof}

A uniformly ready normal policy guarantees long-run optional rate \(r\) if,
on every allowed loss history,
\[
 \liminf_{T\to\infty}\frac1T\int_0^T u(t)\,dt\geq r.
\]
The optimal guaranteed rate is the supremum of these guarantees over
uniformly ready normal policies, allowing paid initialization.

\begin{theorem}[Exact initialized upkeep]
\label{thm:maintenance}
With paid initialization, the maximum guaranteed long-run optional rate under
uniform \(H\)-readiness is \(s-C(H)\) when \(C(H)\leq s\).
If \(C(H)>s\), indefinite uniform readiness is impossible even with no
optional work.
\end{theorem}
\begin{proof}
On maximum feedback put \(Q=\sum_i p_i\). For a uniformly ready policy,
upper reflection and \eqref{eq:normal-budget} imply, almost everywhere,
\[
 Q'\leq \sum_i v_i-\sum_i g_i(p_i)\leq s-u-C(H).
\]
Hence for every duration \(T\),
\begin{equation}
 \int_0^T u(t)\,dt
 \leq [s-C(H)]T+Q(0)-Q(T).
 \label{eq:storage-bound}
\end{equation}
The bounded storage term proves the asymptotic upper bound on this allowed
history. If \(C(H)>s\), the same inequality contradicts \(u\geq0\) for
sufficiently large \(T\).

For attainment, initialize at a minimizing \(p^*\), allocate
\(v_i=g_i(p_i^*)\), and use \(u=s-C(H)\). Under maximum loss the target is
stationary. Under every smaller loss history scalar comparison preserves
\(p\geq p^*\). The ready set is upward closed by the same comparison, so
every request meets \(H\). This constant optional allocation attains the
bound. It does not assert that the actual state is stationary under smaller
losses, or that \(p^*\) is reachable from cold with the normal budget.
\end{proof}

\section{Concentration under unequal proportional losses}
\label{sec:concentration}

For the remainder of the core results let \(g_i(p)=\gamma_i p\), with all
\(\gamma_i>0\), and write \(d_i=\gamma_iM_i\).
A coordinate is \emph{partial} when \(0<p_i<M_i\).
If a serial stage for module \(i\) begins at global time \(t_0\), has capacity
\(b>d_i\), and the module's initial preparation was \(p_i\), its completion
time \(t_1\) satisfies
\begin{equation}
 e^{\gamma_i t_1}
 =\frac{b e^{\gamma_i t_0}-\gamma_i p_i}{b-d_i}.
 \label{eq:proportional-stage}
\end{equation}
The waiting state is \(p_i e^{-\gamma_i t_0}\). A subfull module cannot
complete at capacity \(b\leq d_i\), even if its initial preparation is
arbitrarily close to full.

\begin{theorem}[Unequal-rate concentration]
\label{thm:concentration}
For every \(s\geq0\) and finite \(H\geq0\), every minimizer defining \(C(H)\)
has at most one partial coordinate. Thus a minimizing state consists of a
fully prepared set, at most one partial module, and otherwise cold modules.
\end{theorem}
\begin{proof}
Suppose a minimizing state had at least two partial coordinates. Choose an
attaining serial schedule, transferring all initially full modules at time
zero. Let \(j,k\) be consecutive partial modules in its remaining order.
Every intervening module is cold. Denote the start and end of stage \(j\)
by \(t_0,t_j\), the end of stage \(k\) by \(t_k\), and their capacities by
\(b_j,b_k\). Put
\[
 D_j=b_j-\gamma_jM_j>0,\qquad
 D_k=b_k-\gamma_kM_k>0.
\]
These inequalities follow from actual finite stage feasibility, without any
assumption that the original capacity exceeds every full loss.

The total duration \(L\geq0\) of the intervening cold stages is independent
of their absolute start time and of initial preparations of \(j,k\).
Replace \(p_j\) by \(x\), keeping all other earlier stages fixed. Equation
\eqref{eq:proportional-stage} gives
\[
 t_j(x)=\frac1{\gamma_j}
 \log\frac{b_j e^{\gamma_jt_0}-\gamma_jx}{D_j}.
\]
Choose the replacement \(y(x)\) for \(p_k\) to preserve its original
completion time:
\begin{equation}
 \gamma_k y(x)=
 b_k e^{\gamma_k L}
 \left(\frac{b_j e^{\gamma_jt_0}-\gamma_jx}{D_j}\right)^{\gamma_k/\gamma_j}
 -D_k e^{\gamma_k t_k}.
 \label{eq:pair-compensation}
\end{equation}
At the original point this gives \(y=p_k\). Both coordinates are interior,
so a two-sided open interval of small variations remains inside their
physical bounds. Their stage durations remain positive by continuity
(or directly by \eqref{eq:proportional-stage}). All intervening cold stages
retain their durations. Once stage \(k\) ends, global time and the completed
set agree with the original schedule; every later unserved module has its
unchanged passive state at that same time. The suffix therefore reproduces
the original final deadline.

Only the selected pair's upkeep varies:
\begin{equation}
 c(x)=\gamma_jx+
 b_k e^{\gamma_k L}
 \left(\frac{b_j e^{\gamma_jt_0}-\gamma_jx}{D_j}\right)^{\gamma_k/\gamma_j}
 -D_k e^{\gamma_k t_k}.
 \label{eq:pair-cost}
\end{equation}
If \(\gamma_k<\gamma_j\), this is strictly concave on the open variation
interval. At least one sufficiently small one-sided variation strictly
decreases its value. If \(\gamma_k\geq\gamma_j\), differentiation gives
\begin{equation}
 c'(x)=\gamma_j-
 \frac{b_k\gamma_k}{D_j}
 e^{\gamma_k L}e^{(\gamma_k-\gamma_j)t_j(x)}<0.
 \label{eq:pair-derivative}
\end{equation}
Here \(b_k\geq b_j>D_j\), the times are nonnegative, and
\(\gamma_k\geq\gamma_j>0\). A small increase of \(x\) again strictly lowers
upkeep while \eqref{eq:pair-compensation} keeps \(y\) interior.
Either case contradicts minimality. This proves the assertion for every
minimum, without an iterative exchange or an unproved endpoint limit.
\end{proof}

\begin{corollary}[A restriction on an optimizer's cold prefix]
\label{cor:prefix}
Suppose a minimizing state has a partial module \(k\). In any serial order
for that state that meets deadline \(H\) and transfers its initially full
modules at time zero, every cold module \(j\) before \(k\) satisfies
\(\gamma_j>\gamma_k\).
\end{corollary}
\begin{proof}
All modules between \(j\) and \(k\) are cold by
Theorem~\ref{thm:concentration}. Apply the same variation at \(x=p_j=0\)
and at the interior coordinate \(p_k\). If \(\gamma_j\leq\gamma_k\),
\eqref{eq:pair-derivative} is strictly negative. A small positive \(x\)
keeps \(y(x)\) interior and strictly reduces cost while preserving the
deadline. This contradicts minimality.
\end{proof}

This is a restriction on schedules witnessing an upkeep minimum, not an
ordering rule for arbitrary partially prepared states. In particular, equal
coefficients force an optimizing partial coordinate to precede all cold
modules; unequal coefficients do not. Strict positivity is also substantive:
a zero coefficient would permit cost-free partial coordinates and invalidate
the assertion about every minimum.

\section{An explicit finite readiness frontier}
\label{sec:frontier}

We now eliminate continuous optimization over the initial preparation vector.
All statements in this section allow \(s=0\), zero releases, and inaccessible
initial stages. For a full set \(S\), write \(d(S)=\sum_{j\in S}d_j\).
Define the cold-stage duration
\begin{equation}
 c_i(b)=
 \begin{cases}
 \gamma_i^{-1}\log\!\bigl(b/(b-d_i)\bigr),&b>d_i,\\
 +\infty,&b\leq d_i.
 \end{cases}
 \label{eq:cold-stage}
\end{equation}
Let \(G(S)\) be the exact exit time after immediately transferring full
modules \(S\), with all others cold. The serial theorem gives
\begin{equation}
 G(N)=0,\qquad
 G(S)=\min_{\substack{i\notin S\\b(S)>d_i}}
 \{c_i(b(S))+G(S\cup\{i\})\}.
 \label{eq:cold-tail-dp}
\end{equation}
An empty minimum is infinity. An initially full module transfers without
using \eqref{eq:cold-stage}, even if its putative preparation rate would be
inaccessible.

Fix the initially full set \(S\). For \(T\subseteq N\setminus S\), let
\(L_S(T)\) be the earliest completion time for a cold prefix consisting
precisely of \(T\). Its recurrence is
\begin{equation}
 \begin{split}
 L_S(\varnothing)&=0,\\
 L_S(T)&=\min_{\substack{j\in T\\b(S\cup(T\setminus\{j\}))>d_j}}
 \left\{L_S(T\setminus\{j\})
       +c_j\bigl(b(S\cup(T\setminus\{j\}))\bigr)\right\}.
 \end{split}
 \label{eq:cold-prefix-dp}
\end{equation}
Unreachable prefix states have value infinity. A waiting partial module does
not affect these durations.

Choose disjoint \(S,T\) and \(i\notin S\cup T\). Set
\[
 t=L_S(T),\qquad b=b(S\cup T),\qquad R=G(S\cup T\cup\{i\}).
\]
Consider the candidate only if \(t,R<\infty\) and \(b>d_i\).
If module \(i\) has initial preparation \(q\), its completion after that
prefix occurs at
\begin{equation}
 \Phi_i(t;q,b)=\frac1{\gamma_i}
 \log\frac{b e^{\gamma_i t}-\gamma_i q}{b-d_i}.
 \label{eq:partial-completion}
\end{equation}
For \(0\leq q\leq M_i\), this time is at least \(t\), and
\begin{equation}
 \frac{\partial\Phi_i}{\partial t}
 =\frac{b e^{\gamma_i t}}{b e^{\gamma_i t}-\gamma_i q}>0.
 \label{eq:prefix-monotonicity}
\end{equation}
Thus the earliest prefix is optimal even though the waiting partial
preparation decays. The remaining cold suffix is autonomous, so its best
duration is \(R\).

Writing \([z]_+=\max\{0,z\}\), the least initial preparation for this
displayed strategy to meet \(H\) is
\begin{equation}
 q_{S,T,i}(H)=
 \frac{\left[b e^{\gamma_i t}
       -(b-d_i)e^{\gamma_i(H-R)}\right]_+}{\gamma_i}.
 \label{eq:frontier-candidate}
\end{equation}
Discard candidates with \(q_{S,T,i}(H)>M_i\); the others have cost
\(K_{S,T,i}(H)=d(S)+\gamma_i q_{S,T,i}(H)\).
The endpoint \(q=0\) is cold. The endpoint \(q=M_i\), when the formula is
defined, is a feasible schedule that voluntarily delays a full module.
A full state at \(b\leq d_i\) is represented by adding \(i\) to the initial
full set, never by interpreting a singular partial-stage formula.

\begin{theorem}[Unequal-rate frontier]
\label{thm:frontier}
For every finite \(H\geq0\),
\begin{equation}
 C(H)=\min\left(
 \{d(S):G(S)\leq H\}
 \ \cup\
 \{K_{S,T,i}(H):(S,T,i)\text{ is an eligible candidate}\}
 \right).
 \label{eq:explicit-frontier}
\end{equation}
Eligibility is exactly the finite-prefix, finite-suffix, strict-stage, and
preparation-box conditions above. The formula supplies an attaining initial
state and request schedule. It can be evaluated with
\(O(n3^n)\) real scalar evaluations and comparisons and \(O(2^n)\) working
storage, besides output certificates.
\end{theorem}
\begin{proof}
Every candidate is feasible by its displayed serial construction. Pure
full/cold candidates are feasible by \eqref{eq:cold-tail-dp}; the all-full
set is always one of them. Their costs therefore cannot be below \(C(H)\).

Conversely choose an attained minimum. By
Theorem~\ref{thm:concentration}, it has a full set \(S\) and at most one
partial coordinate. A state without a partial coordinate is already in the
first family. Otherwise let \(i\) be its partial coordinate and choose an
attaining serial order after immediately transferring \(S\). Its other
modules form a cold prefix \(T\) and a cold suffix. Every actual stage is
accessible. Replace the prefix by \eqref{eq:cold-prefix-dp} and the suffix
by \eqref{eq:cold-tail-dp}. Equation \eqref{eq:prefix-monotonicity} shows that
the earlier prefix cannot hurt; the shorter cold suffix cannot hurt either.
Finally reduce initial preparation of \(i\) to
\eqref{eq:frontier-candidate} if necessary. This gives a listed feasible
candidate with no greater cost. Equality follows.

The cold-tail table has \(O(n2^n)\) transitions. Across all full sets \(S\),
there are \(3^n\) disjoint pairs \((S,T)\). The total number of prefix
transitions is \(n3^{n-1}\), and the number of distinguished partial positions
has the same bound. Evaluate one full set's prefix table at a time and retain
only the best candidate, together with the global cold-tail table.
This uses \(O(2^n)\) working storage; an attaining order can be reconstructed
by storing or recomputing the selected prefix certificate.
\end{proof}

The bound counts evaluations and comparisons of real logarithmic and
exponential expressions. It is not a polynomial-time result, a rational-field
bound, or a certified bit-complexity claim. Floating implementations need
explicit comparison tolerances or interval certificates at eligibility
boundaries. Corollary~\ref{cor:prefix} permits pruning prefixes containing
\(\gamma_j\leq\gamma_i\), but the unpruned formula is already exact.
Equation \eqref{eq:explicit-frontier} also gives
\(C(0)=\sum_i d_i\), and \(C(H)=0\) exactly when \(G(\varnothing)\leq H\).

\begin{corollary}[Common coefficient]
\label{cor:common-rate}
If all coefficients equal \(\gamma>0\), only empty cold prefixes are needed
in \eqref{eq:explicit-frontier}. The evaluation count reduces to
\(O(n2^n)\). For rational model parameters and rational
\(X=e^{\gamma H}\), all operations can be performed exactly in the rational
field.
\end{corollary}
\begin{proof}
Corollary~\ref{cor:prefix} excludes a cold predecessor of an optimizing
partial module. Put \(P(S)=e^{\gamma G(S)}\), with infinite values retained.
Then \(P(N)=1\), and accessible transitions obey
\[
 P(S)=\min_{\substack{i\notin S\\b(S)>d_i}}
       \frac{b(S)}{b(S)-d_i}P(S\cup\{i\}).
\]
For an empty-prefix candidate the required preparation becomes
\[
 \frac1\gamma
 \left[b(S)-(b(S)-d_i)\frac{X}{P(S\cup\{i\})}\right]_+.
\]
Keep the same accessibility and box tests and the full/cold candidates
\(P(S)\leq X\). These are rational expressions and comparisons under the
stated arithmetic assumptions. Rational \(H\) alone does not make \(X\)
rational.
\end{proof}

\begin{proposition}[A cold prefix can improve the global optimum]
\label{prop:cold-prefix}
With unequal coefficients, restricting an optimizing partial module to the
first nonfull position can strictly increase the minimum upkeep, already for
two modules with every initial cold stage accessible.
\end{proposition}
\begin{proof}
Take \(s=1\), \(H=\log3\), and
\[
 (M_A,a_A,\gamma_A)=(1/2,3/10,1),\qquad
 (M_B,a_B,\gamma_B)=(1/2,1/10,1/2).
\]
For a state with only the indicated coordinate partial, the four
order/coordinate choices require upkeep
\[
\begin{array}{c|c|c}
 \text{partial coordinate}&\text{order}&\text{required upkeep}\\ \hline
 A&AB&29/1352\\
 B&AB&(26\sqrt2-21\sqrt3)/20\\
 B&BA&1-\tfrac34\sqrt{18/11}\\
 A&BA&7/45.
\end{array}
\]
These follow directly from \eqref{eq:frontier-candidate}. For example,
cold A takes \(\log2\), after which B has capacity \(13/10\);
cold B takes \(2\log(4/3)\), after which A has capacity \(11/10\).
The second cost is positive and rationalizes to
\[
 C_*=\frac{29}{20(26\sqrt2+21\sqrt3)}<\frac{29}{1352}.
\]
For the strict inequality use \(\sqrt2>7/5\), \(\sqrt3>17/10\).
Also \(\sqrt{18/11}<13/10\), so the third cost exceeds
\(1/40>29/1352\); the fourth is larger as well. Every initially full
module costs at least \(1/4\). The all-cold orders miss the deadline because
their exponentiated durations are \(1352/441>3\) and \(88/27>3\).

Concentration and seriality make these comparisons exhaustive. Hence the
global minimum is \(C_*\), attained with
\(p_A=0\), \(p_B=(26\sqrt2-21\sqrt3)/10\), and cold A before partial B.
The best partial-first candidate costs \(29/1352>C_*\).
\end{proof}
